\subsection{Закон вращетельного движения}
%%%%%%%%%%%%%%%%%%%%%%%%%%%%%%%%%%%%%%%%%%%%%%%%%%%%
%%%%%%%%%%%%%%%%%%%%%%%%%%%%%%%%%%%%%%%%%%%%%%%%%%%% 
\z{На концах тонкого однородного стержня длиной $l$ и массой $3m$ прикреплены шарики, массы которых $m$ и $2m$ соответственно. Определить момент инерции такой системы относительно оси, перпендикулярной стержню и проходящей через:
	\lpu шарик массой $m$; \lpuend
	\lpu точку, отстоящую на $r = \frac{1}{3}l$ от шарика массой $m$; \lpuend
	\lpu середину стержня. \lpuend
	}
%
{\puan  $J=3ml^2;$~\puan  $J=\frac{4}{3}ml^2;$~\puan $J=ml^2$}
%%%%%%%%%%%%%%%%%%%%%%%%%%%%%%%%%%%%%%%%%%%%%%%%%%%%	
\z[{\begin{wrapfigure}[5]{r}{0.2\textwidth}%
		\vspace{-35pt}%
		\begin{flushright}%
 			\includegraphics[width=0.2\textwidth]{8.3.2}%
		\end{flushright}%
	\end{wrapfigure}%
	}]%
	{Тонкостенный цилиндр радиуса $R$ раскрутили до угловой скорости $\omega$ и поставили в угол, как показано на рисунке. Коэффициент трения между стенками угла и цилиндром равен $\mu$. Сколько оборотов $N$ цилиндр сделает до полной остановки?}
%
{$n = \frac{\omega^2 R(1 + \mu^2)}{4\pi g\mu(1 + \mu)}$}
%%%%%%%%%%%%%%%%%%%%%%%%%%%%%%%%%%%%%%%%%%%%%%%%%%%% 
%%%%%%%%%%%%%%%%%%%%%%%%%%%%%%%%%%%%%%%%%%%%%%%%%%%%	
\z[{\begin{wrapfigure}[3]{r}{0.2\textwidth}%
		\vspace{-45pt}%
		\begin{flushright}%
 			\includegraphics[width=0.2\textwidth]{8.3.3}%
		\end{flushright}%
	\end{wrapfigure}%
	}]%
	{Определите угловое ускорение $\varepsilon$ блока радиуса $R$ с моментом инерции $J$, вызванное двумя грузами массы $m_1$ и $m_2$, закрепленными на концах нити, перекинутой через блок, если нить не проскальзывает по блоку.}
%
{$\omega = \frac{(m_1 - m_2)gR}{J + m_1 R^2 + m_2 R^2}$}
%%%%%%%%%%%%%%%%%%%%%%%%%%%%%%%%%%%%%%%%%%%%%%%%%%%%
\z[{\begin{wrapfigure}[6]{r}{0.3\textwidth}%
		\vspace{+25pt}%
		\begin{flushright}%
 			\includegraphics[width=0.3\textwidth]{8.3.4}%
		\end{flushright}%
	\end{wrapfigure}%
	}]%
	{На гладком горизонтальном столе находится брусок массы $m_1$. На нем укреплен тонкостенный цилиндр массы $m_2$ и радиуса $R$, который может без трения вращаться вокруг своей оси. На цилиндр намотана невесомая тонкая нить, за конец которой тянут с горизонтальной силой $F$. Найдите ускорение $a$ бруска и угловое ускорение $\alpha$ цилиндра.}
%
{$a = \frac{F}{m_1 + m_2}$, $\omega = \frac{F}{m_2 R}$}
%%%%%%%%%%%%%%%%%%%%%%%%%%%%%%%%%%%%%%%%%%%%%%%%%%%%	
\z[{\begin{wrapfigure}[6]{r}{0.3\textwidth}%
		\vspace{+15pt}%
		\begin{flushright}%
 			\includegraphics[width=0.3\textwidth]{8.3.5}%
		\end{flushright}%
	\end{wrapfigure}%
	}]%
	{Шару сообщили ударом скорость $v_0$ вдоль горизонтальной поверхности стола и вращение вокруг его горизонтального диаметра, перпендикулярного скорости (см. рисунок). В результате скорость шара в течение времени $t_0$ увеличивалась, а затем шар стал двигаться с постоянной скоростью.
	\lpu Найдите эту постоянную скорость $v$.\lpuend
	\lpu На какое расстояние $s$ от места удара удалится шар за время $3t_0$ после удара? \lpuend
	
	\noindent
	Коэффициент трения скольжения между поверхностями шара и стола равен $\mu$.
}
%
{\puan $v = v_0 + \mu g t_0;$~\puan $s = 3v_0 t_0 + \frac{5}{2}\mu g t_0^2$}